
The Euler theory is characterized by the Kirchoff/Love/Euler constraint:

\begin{equation*}
\bm{x}^\prime = \lambda \frameBasis
\end{equation*}
%
An immediate consequence of this is that the components $\gamma_{\beta}$ of the shear strain vector $\shearStrain$ on the cross section directors \(\{\frameBasis_{\FrameYZi}\}\) vanish, leaving only the axial component \(\ShearStrainX\):
%
\begin{equation}
\shearStrain = \left( \|\bm{x}^{\prime}\| - 1 \right) \frameBasis = \epsilon \, \frameBasis
\end{equation}
%
Selecting the material coordinate $\reflenx$ to be normalized by the undeformed arc-length is particularly convenient
and furnishes the representation:

\begin{equation}
\left\|\bm{x}^{\prime}\right\| \equiv \frac{\mathrm{d} \deflenx}{\mathrm{d} \reflenx} \equiv \FrameStretch
\equiv 1 + \ShearStrainX
\end{equation}
%
Note 
    \begin{equation*}
    \frac{\mathrm{d} \FrameRotation}{\mathrm{d} \reflenx} = \FrameRotation \CurveStrain^{\times} 
    \qquad\text{ and }\qquad 
    \frac{\mathrm{d} \FrameRotation}{\mathrm{d} \deflenx} = \FrameRotation \CurveStrainScale^{\times} 
    \end{equation*}
  so
    \begin{equation*}
    \frac{\mathrm{d} \FrameRotation}{\mathrm{d} \reflenx} = \FrameStretch \, \FrameRotation \CurveStrainScale^{\times} = \FrameRotation \CurveStrain^{\times}
    \quad\implies\quad 
    \CurveStrainScale = \FrameStretch^{-1}\CurveStrain
    \end{equation*}
and recall the stretch \(\FrameStretch \triangleq \tfrac{\mathrm{d} \deflenx}{\mathrm{d} \reflenx}\). 
\textcite{dill1992kirchhoffs} makes the remark that \emph{``Since the extension of the axis will ultimately be assumed to be very small,
the distinction between the arc lengths of the deformed and nndeformed axes
of the rod is not carefully maintained by the early writers.''}

\subsubsection{Equilibrium}

  \begin{equation*}
  \begin{aligned}
  \frac{\mathrm{d} \ShearResult}{\mathrm{d} \deflenx}+ \lambda^{-1} \CurveStrain \times \ShearResult + \lambda^{-1} \FrameRotation^{\mathrm{t}}\bm{f}_n = 0,\\
  \quad \frac{\mathrm{d} \CurveResult}{\mathrm{d} \deflenx} + \lambda^{-1} \CurveStrain \times \CurveResult + \FrameBasis \times \ShearResult+ \lambda^{-1} \FrameRotation^{\mathrm{t}}\bm{f}_m=0
  \end{aligned}
  \end{equation*}
  Note that this is in \emph{disagreement} with \cite{simo1991geometricallyexact}, where the notation \(\bm{\Omega}\) is used for our \(\CurveStrain\) and \(\bm{K}\) for our \(\CurveStrainScale\).
When constrained for axial rigidity, it follows that \(\FrameStretch = 1\) and the theory of \textcite{cosserat1908theorie}, and Equations~(10) and (11) from \textcite[387-388]{love1944treatise} is recovered.

Multiplying through by the stretch \(\FrameStretch\) furnishes
\begin{equation}
\begin{aligned}
\ShearResult^{\prime}+\CurveStrain \times \ShearResult + \FrameRotation^{\mathrm{t}}\bm{f}_n =0 \\
\CurveResult^{\prime}+\CurveStrain \times \CurveResult + \FrameStretch \, \FrameBasis \times \ShearResult + \FrameRotation^{\mathrm{t}}\bm{f}_{m}=0
\end{aligned}
\label{eq:exact-euler-MatDefEquil}
\end{equation}
where \((~\cdot~)'\) denotes the derivative with respect to the axial coordinate $\xi$ along the beam reference line, 
\(\ShearResult=\FrameRotation^{\mathrm{t}}\bm{n}\) is the vector of material stress resultant components measured in the deformed beam cross-sectional frame with $V_i=\ShearResult \cdot \frameBasis_i$, 
\(\CurveResult=\FrameRotation^{\mathrm{t}}\bm{m}\) is the vector of stress moment components measured in the deformed beam cross-sectional frame with \(M_i=\CurveResult \cdot \frameBasis_i\), \(\bm{K}\) is the vector of deformed beam curvature components measured in the deformed beam cross-sectional frame with $K_i=\bm{K} \cdot \frameBasis_i, f$ is the vector of applied distributed force components measured in the deformed beam cross-sectional frame with $f_i=\bm{f} \cdot \frameBasis_i$, and $\bm{f}_{m}$ is the vector of applied distributed force components measured in the deformed beam cross-sectional frame with $m_i=\bm{m} \cdot \frameBasis_i$. 

The vector $K=k+\kappa$, where $k$ is the vector of initial twist and curvature components measured in the cross-sectional frame of the undeformed beam, with $k_i=\mathbf{k} \cdot \frameBasis_i$, and $\kappa$ is a measure of elastic twist and bending. 
If the undeformed beam is free of initial twist and curvature, then $k_i=0$.
%
Expanding furnishes
\begin{equation}
\begin{aligned}
N^{\prime}\,-\kappa_1 V_1-\kappa_2 V_2&+\bm{f}_n \cdot \frameBasis\,&=0, \\
V_1^{\prime}+\kappa_1 N\, -    \tau V_2&+\bm{f}_n \cdot \frameBasis_1&=0, \\
V_2^{\prime}+\kappa_2 N\, +    \tau V_1&+\bm{f}_n \cdot \frameBasis_2&=0, \\
T^{\prime}\,+\kappa_1 M_2-\kappa_2 M_1&+m_{\tau}&=0, \\
M_1^{\prime}+\kappa_2 T \,-  \tau\, M_2+ \FrameStretch V_1&+m_1&=0 \\
M_2^{\prime}-\kappa_1 T \,+  \tau\, M_1+ \FrameStretch V_2&+m_2&=0
\end{aligned}
\end{equation}
where 
\begin{equation*}
\bm{x}^{\prime} \times \bm{n} 
= \frac{\mathrm{d} \deflenx}{\mathrm{d} \reflenx} \left(V_1 \frameBasis_2 - V_2 \frameBasis_1\right)
\end{equation*}
which agrees with Equations~(6.6-6.11) of \textcite{dill1992kirchhoffs}.

\paragraph{Condensation}

\begin{equation}
\bm{V} = \FrameBasis\times \frac{1}{\lambda}\left(\CurveResult' + \CurveStrain \times \CurveResult + \CurveLoads - \CurveInert \right)
\label{eq:exact-euler-VofM}
\end{equation}
which is expanded with components % on the TODO basis as
\begin{equation*}
\begin{aligned}
V_1= & \frac{1}{\lambda} \left(- M_2' +\mu_1 T-\mu_3 M_1 + \left(\CurveInert - \bm{c}\right) \cdot \frameBasis_2\right) , \\
V_2= & \frac{1}{\lambda} \left(~~~M_1' +\mu_2 T-\mu_3 M_2 - \left(\CurveInert - \bm{c}\right) \cdot \frameBasis_1\right) .
\end{aligned}
\end{equation*}
with
\begin{equation*}
    \CurveInert = \varrho \bm{J}^{\mathrm{C}} \dot {\bm{\omega}} + \bm{\omega} \times \varrho \bm{J}^{\mathrm{C}} \bm{\omega}
\end{equation*}
Using \cref{eq:exact-euler-VofM} in \cref{eq:exact-euler-MatDefEquil} produces
$$
\begin{aligned}
\left(N \FrameBasis + \FrameBasis\times \frac{1}{\lambda}\left(\CurveResult' + \CurveStrain \times \CurveResult + \CurveLoads - \CurveInert \right)\right)^{\prime}+\bm{K} \times \left(N\FrameBasis + \FrameBasis\times \frac{1}{\lambda}\left(\CurveResult' + \CurveStrain \times \CurveResult + \CurveLoads - \CurveInert \right)\right) + \bm{f}_n =0 \\
{T}^{\prime}+\FrameBasis\cdot \bm{K} \times \CurveResult + f_{\tau}=0
\end{aligned}
$$


$$
\begin{aligned}
\left(N \FrameBasis + \FrameBasis\times \frac{1}{\lambda}\left(\CurveResult' + \CurveStrain \times \CurveResult\right) + \FrameBasis\times \frac{1}{\lambda}\left(\CurveLoads - \CurveInert \right)\right)^{\prime} \\
+\bm{K} \times \left(N\FrameBasis + \FrameBasis\times \frac{1}{\lambda}\left(\CurveResult' + \CurveStrain \times \CurveResult\right) + \FrameBasis\times \frac{1}{\lambda}\left( \CurveLoads - \CurveInert \right)\right) + \bm{f}_n =0 \\
\end{aligned}
$$


$$
\begin{aligned}
\left(\lambda^{-1}~\left[
    - M_2'+\mu_1 T-\mu_3 M_1 - \bm{c} \cdot \frameBasis_2
  \right] \right)'
  -\mu_3 \lambda^{-1} \left[
    ~~M_1'+\mu_2 T-\mu_3 M_2 - \bm{c} \cdot \frameBasis_1
  \right]
  +\mu_2 N 
  +\bm{f} \cdot \frameBasis_1
& = L_1 \\
%
\left(\lambda^{-1}~\left[
    ~~M_1'+\mu_2 T-\mu_3 M_2+\bm{c} \cdot \frameBasis_1
  \right] \right)'
  + \mu_3 \lambda^{-1} \left[
    -M_2'+\mu_1 T-\mu_3 M_1 - \bm{c} \cdot \frameBasis_2
  \right]
  - \mu_1 N 
  + \bm{f} \cdot \frameBasis_2 
& = L_2
\\
\begin{aligned}
N'-\mu_2\left\{\lambda^{-1}~\left[
  -M_2'+\mu_1 T-\mu_3 M_1 + \bm{c} \cdot \frameBasis_2
  \right] \right\}
  + \mu_1 / \lambda \left[
  M_1' + \mu_2 T - \mu_3 M_2 + \bm{c} \cdot \frameBasis_1 
  \right]
% ~~~~~~~~~~~~~
  +\bm{f} \cdot \frameBasis_3
\end{aligned} 
& = L_3\\
T'-\mu_2 M_1+\mu_1 M_2+\bm{c} \cdot \frameBasis_3
=\left[\varrho \bm{J}^{\mathrm{C}} \dot{\bm{\omega}}+\bm{\omega} \times\left(\varrho \bm{J}^{\mathrm{C}} \omega\right)\right] \cdot \frameBasis_3 .
\end{aligned}
$$
where
$$
\begin{aligned}
L_1 &= \varrho A \ddot{\bm{x}} \cdot \frameBasis_1 
     -\mu_3 / \lambda \left(\varrho \bm{J}^{\mathrm{C}} \dot{\bm{\omega}}
     +\left.\omega \times\left(\varrho \bm{J}^{\mathrm{C}} \bm{\omega}\right)\right. \right) \cdot \frameBasis_1 
     - \left[\lambda^{-1} \frameBasis_2 \cdot\left(\varrho \bm{J}^{\mathrm{C}} \dot{\bm{\omega}}\right. + \left.\bm{\omega} \times\left.\varrho \bm{J}^{\mathrm{C}} \bm{\omega}\right.\right)\right]'  \\
%
L_2 &= \varrho A \ddot{\bm{x}} \cdot \frameBasis_2 
      -\mu_3 / \lambda \left(\varrho \bm{J}^{\mathrm{C}} \dot{\bm{\omega}}
      +\left.\omega \times\left(\varrho \bm{J}^{\mathrm{C}} \bm{\omega}\right)\right. \right)  \cdot \frameBasis_2
     + \left(\lambda^{-1} \frameBasis_1 \cdot\left(\varrho \bm{J}^{\mathrm{C}} \dot{\bm{\omega}}\right. + \left.\bm{\omega} \times\left.\varrho \bm{J}^{\mathrm{C}} \bm{\omega}\right.\right)\right)' \\
L_3 &= \varrho A  \ddot{\bm{x}} \cdot \frameBasis_3+\lambda^{-1}\left(\varrho \bm{J}^{\mathrm{C}} \dot{\bm{\omega}}+ \bm{\omega} \times\left(\varrho \bm{J}^{\mathrm{C}} \bm{\omega}\right)\right) \cdot \left(\mu_2 \frameBasis_2+\mu_1 \frameBasis_1\right)
\end{aligned}
$$
which agrees with \cite{lacarbonara2012nonlinear}.

\subsubsection{Remarks}

\begin{marsdenbox}
\cite{boyer2011geometrically}
The rotation $\FrameRotation$ can be expressed in terms of only a twisting angle $\vartheta$ and the axial director $\frameBasis$ which depends only on the tangent position $\bm{x}^{\prime}$ (https://doi.org/10.1002/nme.879):
$$
\FrameRotation = \operatorname{Exp}(\vartheta \frameBasis)\left(\left(\FrameBasis \cdot \frameBasis\right) \mathbf{1}+\frac{1}{1+\FrameBasis \cdot \frameBasis}\left(\FrameBasis \times \frameBasis\right) \otimes\left(\FrameBasis \times \frameBasis\right)+\left(\FrameBasis \times \frameBasis\right)^{\times} \right)
$$
The material curvature \(\bm{\kappa}\) is
$$
\bm{\kappa} = \FrameRotation^{\mathrm{t}}(\alpha^{\prime} \frameBasis + \frameBasis \times \frameBasis^{\prime})
$$
Note
$$
\frameBasis \times \frameBasis^{\prime}=\frac{\bm{x}^{\prime}}{\left\|\bm{x}^{\prime}\right\|} \times\left(\frac{\bm{x}^{\prime}}{\left\|\bm{x}^{\prime}\right\|}\right)^{\prime}
=\left\|\bm{x}^{\prime}\right\|^{-2} \, \bm{x}^{\prime} \times \bm{x}^{\prime \prime}
$$

$$
\delta \alpha' =\delta \theta^{\prime}+\left(\frac{\bm{x}^{\prime} \times \bm{x}^{\prime \prime}}{\left\|\bm{x}^{\prime}\right\|^3}\right) \delta \bm{x}^{\prime}=\delta \theta^{\prime}+\frac{1}{\left\|\bm{x}^{\prime}\right\|}\left(\frameBasis \times \frameBasis^{\prime}\right) \delta \bm{x}^{\prime}
$$
$$
\mathbf{f}_{\text {stat }}=\left[\begin{array}{c}
E A\left(1-\frac{1}{\left\|\bm{x}^{\prime}\right\|}\right) \bm{x}^{\prime}+E I_a\left(2 \frac{\left(\bm{x}^{\prime} \cdot \bm{x}^{\prime \prime}\right)^2}{\left\|\bm{x}^{\prime}\right\|^6}-\frac{\left(\bm{x}^{\prime} \cdot \bm{x}^{\prime \prime}\right) \bm{x}^{\prime \prime}+\left\|\bm{x}^{\prime \prime}\right\|^2 \bm{x}^{\prime}}{\left\|\bm{x}^{\prime}\right\|^4}\right) \\
E I_a\left(\frac{\bm{x}^{\prime \prime}}{\bm{x}^{\prime 2}}-\frac{\left(\bm{x}^{\prime} \cdot \bm{x}^{\prime \prime}\right)}{\left\|\bm{x}^{\prime}\right\|^4} \bm{x}^{\prime}\right) \\
0
\end{array}\right]+\left[\begin{array}{c}
G I_p \alpha ' \left(\frac{\bm{x}^{\prime} \times \bm{x}^{\prime \prime}}{\left\|\bm{x}^{\prime}\right\|^3}\right) \\
\mathbf{0}_{3 \times 1} \\
G I_p \alpha'
\end{array}\right]
$$
\end{marsdenbox}

%
By virtue of the constraints $\gamma_1=0=\gamma_2$, the gradient of $\bm{x}$ is $\bm{x}' = \lambda \mathbf{b}_3$ with 
$$
\bm{x}' = u_1'\, \mathbf{e}_1+ u_2'\, \mathbf{e}_2+ \left(1+ u_3'\right)\, \mathbf{e}_3
$$
Hence, by using 
\begin{equation*}
\frameBasis = \FrameRotation \FrameBasis
=\sin \theta_2 \FrameBasis_1-\cos \theta_2 \sin \theta_1 \FrameBasis_2+\cos \theta_1 \cos \theta_2 \FrameBasis
\end{equation*}
and
$$
\begin{aligned}
u_1' &= \sec \theta_1 \tan \theta_2\left(1+ u_3' \right), \\
u_2' &=-\tan \theta_1 \left(1+ u_3'\right)
\end{aligned}
$$
the following kinematic relationships are obtained:
$$
\begin{aligned}
  u_1' & =  \lambda  \sin \theta_2, \\
  u_2' & = -\lambda  \cos \theta_2 \sin \theta_1, \\
1+u_{\FrameXi}' & =  \lambda  \cos \theta_1 \cos \theta_2
\end{aligned}
$$
The latter can be solved to yield the rotation angles $\theta_1$ and $\theta_2$ and the stretch \(\lambda\) as
$$
\begin{aligned}
\theta_1 & =-\arctan \left(\frac{u_2'}{1+u_3'}\right), \\
\theta_2&=~\arctan \left[\frac{u_1'}{\sqrt{\left(u_2'\right)^2+\left(1+ u_3'\right)^2}}\right] \\
\lambda & =\sqrt{\left(u_1'\right)^2+\left(u_2'\right)^2+\left(1+u_{\FrameXi}'\right)^2}
\end{aligned}
$$

\begin{enumerate}[label=(\roman*)]
\item With warping torsion one has:
    \begin{equation*}
    \begin{array}{r}
    N^{\prime}-V_2 \kappa_3+V_3 \kappa_2=0 \\
    V_2^{\prime}-V_3 \kappa_1+N \kappa_3=0 \\
    V_3^{\prime} - \kappa_2 N + V_2 \kappa_1=0 \\
    (M_1 - \Bimoment')^{\prime} - M_2 \kappa_3 + M_3 \kappa_2 = 0 \\
    M_2^{\prime}-M_3 \kappa_1 + (M_1- \Bimoment^{\prime}) \kappa_3- \FrameStretch V_3 = 0 \\
    M_3^{\prime} - (M_1 - \Bimoment) \kappa_2+M_2 \kappa_1 + \FrameStretch V_2 = 0
    \end{array}
    \end{equation*}

\item Neglecting shear, one has \parencite{cosserat1908theorie,ericksen1957exact}
  $$
  \begin{aligned}
  \frac{\mathrm{d} \bm{n}}{\mathrm{d} \deflenx}+ \FrameStretch^{-1} \bm{f}_n = 0 \\
  \quad \frac{\mathrm{d} \bm{m}}{\mathrm{d} \deflenx} + \frameBasis \times \bm{n} + \FrameStretch^{-1} \bm{f}_m = 0
  \end{aligned}
  $$
  with components:
  $$
  \begin{aligned}
  \frac{\mathrm{d} M_1}{\mathrm{d} \deflenx}-w_1^{\mathfrak{b}} M_{\mathfrak{b}}+ \bm{f}_m \cdot \frameBasis_1 &=0, \\
  \frac{\mathrm{d} M_2}{\mathrm{d} \deflenx}-w_2^{\mathfrak{b}} M_{\mathfrak{b}}-V^3+ \bm{f}_m \cdot \frameBasis_2&=0, \\
  \frac{\mathrm{d} M_3}{\mathrm{d} \deflenx}-w_3^{\mathfrak{b}} M_{\mathfrak{b}}+V^2+ \bm{f}_m \cdot \frameBasis_3&=0 .
  \end{aligned}
  $$
  Note that this form differs only by a scaling by the stretch $\lambda$

\item Neglecting axial extension, one has Kirchhoff's classical equations (see, eg, \cite{reissner1973onedimensional,dill1992kirchhoffs}):
  $$
  \begin{aligned}
  N^{\prime}\,-\kappa_1 V_1-\kappa_2 V_2&+\bm{f}_n \cdot \frameBasis\,&=0, \\
  V_1^{\prime}+\kappa_1 N\,-    \tau V_2&+\bm{f}_n \cdot \frameBasis_1&=0, \\
  V_2^{\prime}+\kappa_2 N\,+    \tau V_1&+\bm{f}_n \cdot \frameBasis_2&=0, \\
  T^{\prime}\,+\kappa_1 M_2-\kappa_2 M_1&+m_{\tau}&=0, \\
  M_1^{\prime}+\kappa_2 T\,-  \tau\, M_2+ V_1&+m_1&=0 \\
  M_2^{\prime}-\kappa_1 T\,+  \tau\, M_1+ V_2&+m_2&=0
  \end{aligned}
  $$

\item Other

$$
(1+\epsilon)^2=\bm{x}^{\prime}\cdot\bm{x}^{\prime} =  (u_y^{\prime})^2 + (u_z^{\prime})^2 + \left(1+u_{\FrameXi}^{\prime}\right)^2
$$

\begin{equation}
\bm{\gamma}
=\begin{pmatrix}
\sqrt{(u_y^{\prime})^2+(u_z^{\prime})^2+\left(1+u_x^{\prime}\right)^2} - 1 \\
0 \\
0 \\
\end{pmatrix}
\end{equation}

$$
r(\bm{x})=\sqrt{\bm{x} \cdot \bm{x}}
$$


$$
\begin{aligned}
D\, r(\bm{x}) \cdot \bm{\eta} & =\frac{d}{d \varepsilon}[r(\bm{x}+\varepsilon \bm{\eta})]_{\varepsilon=0} \\
& =\frac{d}{d \varepsilon}\left[((\bm{x}+\varepsilon \bm{\eta}) \cdot(\bm{x}+\varepsilon \bm{\eta}))^{\frac{1}{2}}\right]_{\varepsilon=0} \\
& =\left[\frac{1}{2}((\bm{x}+\varepsilon \bm{\eta}) \cdot(\bm{x}+\varepsilon \bm{\eta}))^{-\frac{1}{2}}(2(\bm{x}+\varepsilon \bm{\eta}) \cdot \bm{\eta})\right]_{\varepsilon=0} \\
& =\frac{\bm{x} \cdot \bm{\eta}}{\sqrt{\bm{x} \cdot \bm{x}}}=\frac{\bm{x}}{r(\bm{x})} \cdot \bm{\eta}
\end{aligned}
$$

\end{enumerate}

\begin{ambox}[\gls{frame:exact} \gls{frame:euler} Frame]{box:exact-euler}

\paragraph{Equilibrium}
\begin{equation*}
-\EquilibriumDerivative \bm{A}_* \bm{s}_{\mathrm{p}} = 
\left\{
\begin{aligned}
\ShearResult^{\prime} + \CurveStrain \times \ShearResult \\
\CurveResult^{\prime} + \CurveStrain \times \CurveResult + \FrameStretch \, \FrameBasis \times \ShearResult 
\end{aligned}
\right.
\end{equation*}

\paragraph{Kinematics}
\begin{equation*}
\ShearStrain = \left( \|\bm{x}^{\prime}\| - 1 \right) \FrameBasis
\end{equation*}

\end{ambox}
